\subsection{Configurazione e algoritmo associato}
Alla luce di quanto fatto prima, sembrerebbe naturale chiamare algoritmica qualsiasi procedura costruita a partire da un
automa finito.
In particolare è chiaro che una tale procedura può essere descritta in un linguaggio di
programmazione e implementata in un qualsiasi sistema informatico. L’affermazione reciproca al contrario
sembrerebbe troppo restrittiva: non sembra difendibile che ogni procedura che si voglia definire algoritmica
debba essere realizzata mediante il modello di algoritmi associato agli automi finiti (più avanti vederemo degli esempi). 
Abolire completamente il carattere finito dell’insieme degli stati
svuoterebbe di significato la nozione stessa di algoritmo. Un algoritmo è, per definizione, una procedura descritta da un insieme finito di istruzioni/regole — è proprio questo che lo rende "eseguibile", trascrivibile, comunicabile, implementabile in un programma. Se l'insieme degli stati fosse infinito, la "macchina" non sarebbe più descrivibile con una quantità finita di dati.
Per conservare il carattere finitario auspicato senza rinunciare a fondamentali
capacità di calcolo, si può prescindere dal principio che l’automa non modifica la parola iniziale, dunque nel
leggere un certo simbolo dalla memoria potrà successivamente utilizzare questa informazione per modificare il
proprio stato ma anche per modificare la memoria scrivendo sulla stessa un nuovo carattere.
Questa variante
risulta interessante solo se a sua volta lo scorrimento automatico dei caratteri in lettura viene modificato per
permettere alla macchina di riesaminare i caratteri che sono stati modificati nei passi di calcolo precedenti,
altrimenti la modifica che permette di scrivere non avrebbe alcun impatto sul calcolo effettuato.
\begin{exam}
Consideriamo lgli insiemi
\[
X_0=
\{
0^n1^m \, | \, n,m \in \mathbb N
\},
\quad
X_1=
\{
0^n1^n \, | \, n \in \mathbb N
\} 
\]
Sebbene $X_0$ sia un insieme DFA-decidibile,
$X_1$ non lo è.
Infatti
gli automi non sono in grando di "contare",
quindi non si potrebbe contare il numero di occorenze di
$1$ e $0$
(forniremo una dimostrazione più avanti).
\end{exam}
\begin{figure}[H]
\centering
\begin{tikzpicture}[
    >=stealth', shorten >=1pt, auto,
    every state/.style={circle, draw, fill=white, minimum size=0.8cm, inner sep=1pt}
]
% Stati
\node[state, accepting] (q0) at (0,0)    {$q_0$};
\node[state, accepting] (q1) at (3.5,0)  {$q_1$};
\node[state]            (qp) at (1.75,-2.5) {$q_p$};

% Freccia d'ingresso
\draw[->] ([shift={(-0.7,0.7)}]q0.center) -- (q0);

% Transizioni
\draw[->] (q0) -- node {$1$} (q1);
\draw[->] (q1) -- node {$0$} (qp);

% Cappi
\path[->] (q0) edge[loop above] node {$0$}    (q0);
\path[->] (q1) edge[loop above] node {$1$}    (q1);
\path[->] (qp) edge[loop below] node {$0,1$}  (qp);
\end{tikzpicture}
\caption{DFA che decide il linguaggio $\{0^n1^m \mid n,m\in\mathbb N\}$.}
\label{fig:dfa-0n1m}
\end{figure}
In questa sezione vogliamo introdurre un modello
di calcolo più forte degli automi.
\begin{defi} Una \emph{configurazione ad un nastro di alfabeto $A$ e insieme degli stati $Q$} è una tripla $(f,p,q)$, dove
\begin{enumerate}
\item il \textit{nastro} è una funzione $f : \mathbb{Z} \to A'$ dove $A' := A \cup \{\square\}$ (a volte indicato anche con $\tilde{A}$) e tale che $|f^{-1}(A)| < \infty$,
\item $p \in \mathbb{Z}$,
\item $q \in Q'$ dove $Q' = Q \cup \{q_F^0, q_F^1\}$ essendo $q_F^0$ e $q_F^1$ due stati che non appartengono a $Q$.
\end{enumerate}

L'insieme delle configurazioni per un nastro di alfabeto $A$ e insieme degli stati $Q$ sarà denotato con
\[
Conf_1^{A,Q} := \{(f,p,q) \text{ che verificano i punti {(1)}, {(2)}, e {(3)}}\}
\]

\end{defi}
Il criterio di arresto, che questa volta non è determinato dalla fine della lettura del nastro, deve essere
specificato esplicitamente utilizzando un certo numero di stati terminali, che nella definizione sopra sono $q_F^0$ e $q_F^1$.
Bisogna inoltre fissare uno stato particolare da utilizzare per la codifica della configurazione iniziale a partire dalla parola di ingresso. Questo stato iniziale indicato da $q_0 \in  Q$ determina allora la configurazione
iniziale tramite una funzione di input $\epsilon()$.
Lo stato iniziale sarà d’ora in poi indicato da $q_0$.
\begin{defi}[Machcina di Turing]
     Dato un alfabeto $A$ e un insieme finito di stati $Q$, la \emph{macchina di Turing di alfabeto $A$ e stati $Q$} è un'applicazione
\[
\mu : Q \times \tilde{A} \to Q' \times \tilde{A} \times \{+1,-1\}.
\]

\end{defi}

In realtà  vorremmo che la macchina di Turing sia allo stesso tempo i seguenti tre oggetti:
\[
Q\times A \to A, \qquad Q\times A \to \{-1,+1\}, \qquad Q \times A \to Q
\]
Le precedenti funzioni rappresentano rispettivamente il sovrascrimento di un carattere, il cambio di stato e lo spostamento del puntatore.
La funzione $\mu$ della definizione precedente formalizza in un'unica funzione i tre precedenti dati, infatti
\[
\mu : Q \times \tilde{A} \to Q' \times \tilde{A} \times \{+1,-1\}
\]
prende in input una coppia $(q,a)$, dove:
\begin{itemize}
\item $q \in Q$ è lo \textit{stato corrente} della macchina,
\item $a \in \tilde{A}$ è il \textit{carattere letto} sotto il puntatore (può essere un simbolo di $A$ oppure il blank $\square$),
\end{itemize}
e restituisce una tripla $(q', a', m)$, dove:
\begin{itemize}
\item $q' \in Q'$ è il \textit{nuovo stato} in cui passa la macchina (può essere uno stato ``di lavoro'' in $Q$, oppure uno dei due stati terminali $q_F^0, q_F^1$ se il calcolo deve fermarsi),
\item $a' \in \tilde{A}$ è il \textit{carattere da scrivere} al posto di $a$ nella cella corrente,
\item $m \in \{+1,-1\}$ è lo \textit{spostamento} del puntatore: $+1$ a destra, $-1$ a sinistra.
\end{itemize}
In altre parole, $\mu$ è la \emph{tabella delle regole}: ``se sono nello stato $q$ e leggo il carattere $a$, allora scrivo $a'$, mi sposto di $m$ e passo allo stato $q'$''. 

\begin{defi}
Sia $\mu$ una macchina di Turing di alfabeto $A$ e insieme degli stati $Q$, l'\emph{algoritmo formale} $\mathcal{A}_\mu$ associato a $\mu$ è la tupla $(A,\{0,1\}, Conf_1^{A,Q},\epsilon, \tau, \delta)$: dove
\begin{enumerate}
\item l'insieme delle configurazioni è dato da
\[
Conf_1^{A,Q}
\]

\item la funzione di input $\epsilon : A^* \to Conf_1^{A,Q}$ dove $\epsilon(w) = (f_w, 1, q_0)$ con $f_w$ corrisponde alla funzione nastro per cui
\[
f_w(n) = \begin{cases} w(n) & \text{se } 1 \le n \le |w| \\ \square & \text{altrimenti} \end{cases}
\]

\item la funzione di transizione
\[
\tau : Conf_1^{A,Q} \to Conf_1^{A,Q}
\]
è definita come $\tau((f,p,q)) = (f',p',q')$ dove supposta $\mu(q,f(p)) = (q',a',m)$ si porrà $p' = p+m$ e
\[
f'(x) = \begin{cases} f(x) & \text{se } x \neq p \\ a' & \text{se } x = p. \end{cases}
\]

\item la funzione di uscita $\delta$ associa ad una configurazione data $(f,p,q)$ dove $q \in \{q_F^0, q_F^1\}$ il valore $0$ se $q = q_F^0$ e il valore $1$ se $q = q_F^1$.
\end{enumerate}
\end{defi}

Cerchiamo di capire come si esegue un passo di calcolo in questo caso.
\begin{rema}
Data una configurazione $(f,p,q) \in Conf_1^{A,Q}$ (nastro, posizione del puntatore, stato corrente), un passo di calcolo si esegue così:
\begin{enumerate}
\item si legge il carattere sotto il puntatore: $a := f(p)$;
\item si applica $\mu$: $\mu(q,a) = (q',a',m)$;
\item si aggiorna lo stato: il nuovo stato è $q'$;
\item si aggiorna la posizione: $p' = p+m$ (il puntatore si sposta di una cella a destra o sinistra);
\item si aggiorna il nastro: resta identico ovunque tranne che nella cella $p$, dove viene scritto $a'$:
\[
f'(x) = \begin{cases} f(x) & x \neq p \\ a' & x = p. \end{cases}
\]
\end{enumerate}
Questo definisce la funzione di transizione $\tau : Conf_1^{A,Q}\to Conf_1^{A,Q}$, $\tau(f,p,q)=(f',p',q')$.

\end{rema}
Vediamo ora l'avvio e l'arresto del calcolo.
\begin{rema}\
\begin{itemize}
\item \emph{Avvio:} data una parola di input $w \in A^*$, la configurazione iniziale è $\epsilon(w) = (f_w, 1, q_0)$: il nastro contiene $w$ scritta a partire dalla cella $1$ (blank altrove), il puntatore parte dalla cella $1$ (inizio della parola), lo stato iniziale è $q_0$.
\item \emph{Iterazione:} si applica ripetutamente $\tau$ per ottenere la sequenza di configurazioni $$\epsilon(w), \tau(\epsilon(w)), \tau^2(\epsilon(w)), \dots$$  questa è l'esecuzione del calcolo, passo dopo passo.
\item \emph{Arresto:} il calcolo si ferma quando si raggiunge una configurazione il cui stato è $q_F^0$ oppure $q_F^1$ (questi sono gli unici stati per cui $\mu$ non è più definita/rilevante, essendo terminali).
\item \emph{Output:} se ci si ferma in $q_F^1$ l'output è $1$ (accetta); se ci si ferma in $q_F^0$ l'output è $0$ (rifiuta). Se non ci si ferma mai, l'algoritmo non produce output su quell'input (il calcolo diverge).
\end{itemize}
In sintesi, $\mu$ è la ``regola locale'' (cosa fare dato stato+carattere letto); $\tau$ applica questa regola a una configurazione globale producendone la successiva; $\epsilon$ dice come tradurre l'input iniziale in una configurazione; l'iterazione di $\tau$ a partire da $\epsilon(w)$, fino a un eventuale stato terminale, è \emph{l'esecuzione dell'algoritmo} $\mathcal{A}_\mu$ sull'input $w$.
\end{rema}

\subsection{Decidibilità per macchina di Turing}
\begin{defi}
Sia $X \subset A^*$, allora $X$ si dice \emph{decidibile tramite Macchina di Turing} (o \emph{Turing-decidibile}) se esiste $Q$ insieme finito di stati e una macchina di Turing
di alfabeto $A$ e insieme degli stati $Q$ per cui l'algoritmo associato $\mathcal{A}_\mu$ decide l'insieme $X$.

\end{defi}

Analogamente si ottiene la definizione di insieme \emph{semi-decidibile} per macchina di Turing se esiste una macchina di Turing tale che l'algoritmo associato semi-decide $X$.

\begin{prop}
Ogni insieme decidibile per macchina di Turing è anche semidecidibile per macchina di
Turing.
\begin{proof}
    Per dimostrare questa proposizione è sufficiente costruire una macchina 
$\tilde \mu$ che per ogni
parola $w\in A^*$ e $w \notin X$ entri in una configurazione non terminale.
A questo scopo, sia $\mu$ la macchina che decide $X$ e sia $Q$ il suo insieme degli stati, andiamo a modificarla
nelle transizioni che coinvolgono il suo stato finale non accettante.
Consideriamo la macchina $\mu^{*}$ di alfabeto $A$ e insieme degli stati
$Q^{*}=Q\cup\{q^{*}\}$ che in corrispondenza delle transizioni di $\mu$
che coinvolgono lo stato non accettante opera una transizione verso una
configurazione $(f,p,q^{*})$ per cui
\[
\tau(f,p,q^{*})=(f,p',q^{*})
\]
che porta la macchina in una computazione che non termina.
Più precisamente, data la macchina
\[
\mu : Q\times A' \longrightarrow Q' \times A' \times \{-1,+1\},
\]
si costruisce la macchina
\[
\mu^{*}:Q^{*}\times A'
\longrightarrow
Q^{*'}\times A'\times\{-1,+1\},
\]
dove
\[
Q^{*}=Q\cup\{q^{*}\}
\]
e
\[
\mu^{*}(q,a)=
\begin{cases}
\mu(q,a)
&
\text{se }\mu(q,a)=(q',a',m)
\text{ e }q'\neq q_{F}^{0},
\\[1ex]
(q^{*},a,+1)
&
\text{altrimenti.}
\\

\end{cases}
\]
\end{proof}
\end{prop}
\begin{rema} 
    
    Il modello di calcolo ``macchina di Turing'' è più potente del modello ``automa a stati finiti'', in quanto per ogni automa $(T, q_0, F)$ resta definita una macchina di Turing $\mu$ sullo stesso alfabeto e sullo stesso insieme di stati per cui la computazione dell'algoritmo associato alla macchina $\mu$ coincide con l'algoritmo associato all'automa, più precisamente è sufficiente prendere
\[
\mu(q,a) :=
\begin{cases}
(T(q,a), a, +1) & a \ne \square \\[4pt]
(q_F^1, \square, +1) & \text{se } q \in F \\[4pt]
(q_F^0, \square, +1) & \text{se } q \notin F
\end{cases}
\]
\end{rema}
\begin{rema}
    Come nel caso degli automi
    possiamo associare ad ogni macchina
    un grafo etichettato che 
    ha come vertici gli stati 
    e le etichette sono 
    della forma
    $a/b, m$ e indicano che
    se si legge $a$ si sostituisce con $b$
    e si incrementa il puntatore di $m$.
\end{rema}
\begin{exam}
    \label{macchina turing 0n1n}
    Vogliamo mostrare che l'insieme
    \[
    X = \{0^n1^n \, | \, n \in \mathbb N\}
    \]
    è decidibile secondo macchina di Turing.
    Questo si può fare
    mediante una serie di cancellazioni successive della coppia di caratteri costituita dallo zero più a destra
    e dell’1 più a sinistra.
\end{exam}
\begin{figure}[H]
\centering
\begin{tikzpicture}[
    scale=1, transform shape,
    >=stealth', shorten >=1pt, auto,
    every state/.style={circle, draw, fill=white, minimum size=0.8cm, inner sep=1pt}
]
% Stati della riga principale
\node[state] (q0) at (0,0)    {$q_0$};
\node[state] (q1) at (3,0)    {$q_1$};
\node[state] (q2) at (6,0)    {$q_2$};
\node[state] (q3) at (9,0)    {$q_3$};
\node[state] (q4) at (12,0)   {$q_4$};

% Stati finali
\node[state, accepting] (qF0) at (7.5,3.8)  {$q_F^0$};
\node[state, accepting] (qF1) at (0,-3)     {$q_F^1$};

% Freccia d'ingresso
\draw[->] ([shift={(-0.7,0.7)}]q0.center) -- (q0);

% Transizioni lungo la riga principale
\draw[->] (q0) -- node {$0/\square,+1$} (q1);
\draw[->] (q1) -- node {$1/1,+1$} (q2);
\draw[->] (q2) -- node {$\square/\square,-1$} (q3);
\draw[->] (q3) -- node {$1/\square,-1$} (q4);

% Cappi
\draw[->] (q1) to[loop above, looseness=8] node {$0/0,+1$} (q1);
\draw[->] (q2) to[loop below, looseness=8] node {$1/1,+1$} (q2);
\draw[->] (q4) to[loop above, looseness=8] node {$1/1,-1$} (q4);
\draw[->] (q4) to[loop below, looseness=8] node {$0/0,-1$} (q4);

% Transizioni verso q_F^0 (rifiuto)
\draw[->] (q0) to[bend left=30]  node[pos=0.4, above] {$1/$}       (qF0);
\draw[->] (q1) to[bend left=20]  node[pos=0.5, above] {$\square/$} (qF0);
\draw[->] (q2) to[bend left=20]  node[pos=0.5, above] {$0/$}       (qF0);
\draw[->] (q3) to[bend left=25]  node[pos=0.5, left]  {$0/$}       (qF0);
\draw[->] (q3) to[bend right=25] node[pos=0.5, right] {$\square/$} (qF0);

% Transizione verso q_F^1 (accettazione)
\draw[->] (q0) -- node[swap] {$\square/$} (qF1);

% Ritorno da q_4 a q_0
\draw[->] (q4) to[bend left=30] node[pos=0.5, below] {$\square/\square,+1$} (q0);
\end{tikzpicture}
\caption{Macchina di Turing che decide $X=\{0^n1^n \mid n\geq 0\}$.}
\label{fig:turing-0n1n}
\end{figure}

